\documentclass[11pt]{article}
\usepackage[margin=1in]{geometry}
\usepackage{enumitem}
\usepackage{booktabs}
\title{Efficient Model Adaptation Program Write-up}
\author{ }
\date{March 2026}

\begin{document}
\maketitle

\section{Mission}
This repository demonstrates how parameter-efficient fine-tuning (PEFT),
retrieval-augmented generation (RAG), and quantized inference can be
combined to ship grounded assistants that run on accessible hardware.
Each subproject highlights a different slice of that story: a technical
VQA agent and a medical imaging study that quantifies LoRA/BitFit
trade-offs.

\section{Completed Work}
\subsection{PEFT-Thyroid LLaVA-Med Study}
\begin{itemize}[leftmargin=*]
\item Curated and cleaned the DDTI thyroid ultrasound dataset (hosted on
    Kaggle) plus an updated derivative tailored for PEFT experiments.
\item Adapted LLaVA-Med using LoRA and BitFit, releasing public
    checkpoints and notebooks that walk through dataset prep, baseline
    establishment, hyperparameter sweeps, and evaluation.
\item Captured reproducibility guidance in
    \texttt{src/peft-thyroid/README.md}, including how to attach
    datasets/weights inside Kaggle notebooks to work around local VRAM
    limits.
\end{itemize}

\subsection{Strategic Planning for PEFT-VQA}
\begin{itemize}[leftmargin=*]
\item Captured user personas, evaluation metrics, and a 12-task execution
    plan directly inside this write-up so every agent has the same brief.
\item Defined the agent architecture (retrieval layer, tool adapters,
    PEFT-tuned reasoning model, quantization + evaluation loop) and the
    communication plan for stakeholders.
\end{itemize}

\subsection{MLX Voice Assistant Prototype}
\begin{itemize}[leftmargin=*]
\item Built a fully local speech loop that records microphone audio,
    transcribes with Parakeet, queries a small Ollama model, and vocalizes
    the cleaned reply through Kokoro running on Apple MLX.
\item Captured the operational playbook in
    \texttt{app/voice-ai/readme.md}: `uv sync`, ensure Homebrew installs
    (`espeak-ng`, `portaudio`), and launch via
    \verb|ESPEAK_DATA_PATH="$(brew --prefix espeak-ng)/share/espeak-ng-data"|
    followed by \texttt{uv run src/agent.py}.
\item Documented customization knobs (\texttt{LOCAL\_STT\_ID},
    \texttt{LOCAL\_TTS\_VOICE}, \texttt{OLLAMA\_MODEL}) so future demos
    can swap voices/models without code edits.
\end{itemize}

\section{In-Flight and Planned Work}
\subsection{Retrieval-Centric Technical VQA Assistant}
\begin{itemize}[leftmargin=*]
\item \textbf{Personas \\ knowledge base}: finalize prioritized
    workflows (students, engineers, researchers) and ingest high-signal
    CV/ML/graphics sources with robust metadata.
\item \textbf{Retrieval pipeline}: stand up chunking, embeddings, and
    recall/precision harness using lightweight scripts under
    \texttt{src/peft-vqa} plus the shared prompt manifest in
    \texttt{test.json}.
\item \textbf{Tool adapters}: expose code execution, math helpers, PDF
    parsing, and repo search with guardrails so the assistant behaves like
    a true agent.
\item \textbf{PEFT dataset + experiments}: generate instruction-style QA
    pairs, then compare LoRA, QLoRA, and prompt-tuning adapters for
    accuracy, hallucination, and token efficiency.
\item \textbf{Quantization + deployment}: profile 4/8-bit variants,
    measure memory/latency deltas, and route easy vs. hard queries to the
    right model footprint.
\item \textbf{Evaluation loop}: track hallucination/citation checks
    inside \texttt{test.json} + README tables, then promote promising
    metrics back into this write-up for reliability reviews.
\item \textbf{UI + ops}: ship a Streamlit/Gradio/chat UI with upload
    support, telemetry, and containerized deployment scripts.
\end{itemize}

\subsection{PEFT-Thyroid Follow-ups}
\begin{itemize}[leftmargin=*]
\item Explore additional adapters (e.g., IA3) and distillation
    strategies to shrink deployment footprint without sacrificing
    diagnostic fidelity.
\item Integrate dataset alignment checks and model monitoring so
    clinical collaborators can trust distributional stability over time.
\end{itemize}

\section{Roadmap Snapshot}
Table~\ref{tab:roadmap} distills the open tasks currently tracked inside
this write-up. Status values will remain synchronized with that source
of truth.

\begin{table}[h]
    \centering
    \begin{tabular}{@{}ll@{}}
        \toprule
        \textbf{Workstream} & \textbf{Upcoming Deliverables} \\
        \midrule
        Personas + scope & Persona briefs, prioritized use cases \\
        Knowledge base & Source inventory, ingestion spec, embeddings \\
        Retrieval pipeline & Vector service, recall report \\
        Tool adapters & Guarded tool catalog + tests \\
        Baseline model & Agent notebook, telemetry logs \\
        PEFT dataset & Training corpus, taxonomy \\
        PEFT sweeps & Experiment tracker, best checkpoint \\
        Quantization & Benchmark sheet, quantized artifacts \\
        Evaluation loop & Automated factual/grounding eval suite \\
        UI & Chat interface with citations and uploads \\
        Deployment & Docker/infra scripts, ops doc \\
        Story assets & README refresh, demo script, metrics summary \\
        \bottomrule
    \end{tabular}
    \caption{Rolling plan for the technical VQA assistant track.}
    \label{tab:roadmap}
\end{table}

\section{Call to Action}
Short term, the priority is to operationalize the VQA agent foundation:
lock personas, build the retrieval stack, and stand up the baseline
reasoning model so PEFT + quantization experiments have a stable target.
In parallel, continue polishing the MLX voice runtime so it can front
the same reasoning stack once retrieval is ready. Every change should be
reflected in this write-up and surfaced through the repo-orientation
skill so future agents can align instantly.

\end{document}
